% F14b fig:ridders-cost (Chapter 6, margin). Histogram of the evaluations of fn per call over the random study
% (DS9, data/ridders-coverage-N1.csv and -N2.csv): N = 1 (dark, even counts: 2 per level) and N = 2 (muted, odd
% counts: 2 per level plus fn(x)). Unit bins centred on the integers, so the two series interleave. The tail of
% N = 2 at 19 and 21 is the slower convergence of several functions; the 11 restarts of log and sqrt below x = 0.25 are a small part of the last bar (R1-decision T7).
\begin{tikzpicture}
\begin{axis}[eqtufte, width=38mm, height=30mm,
  xmin=5.5, xmax=22.5, ymin=0, ymax=2100,
  xtick={6,10,14,18,22}, ytick={0,1000,2000}, scaled y ticks=false,
  xlabel={evaluations of \code{fn}}, ylabel={cases}]
  \addplot[hist={bins=17, data min=5.5, data max=22.5}, fill=black!30, draw=none]
    table[col sep=comma, comment chars={\#}, y=evaluations] {data/ridders-coverage-N2.csv};
  \addplot[hist={bins=17, data min=5.5, data max=22.5}, fill=black!75, draw=none]
    table[col sep=comma, comment chars={\#}, y=evaluations] {data/ridders-coverage-N1.csv};
  \node[font=\tiny, anchor=south, inner sep=1pt] at (axis cs:14.5, 1900) {$\dorder = 1$};
  \node[font=\tiny, anchor=south, inner sep=1pt, text=black!55] at (axis cs:17.5, 1580) {$\dorder = 2$};
\end{axis}
\end{tikzpicture}
